\subsection{All even terminal lists with at least two unequal pairs at a unit root}\label{proof:R2-C029}
Let $a_1<b_1<\cdots<a_j<b_j$ with $j\ge2$. Put $F=\prod_i h_{a_i}f_{b_i}$ and $H=pT^LF$, where $L\ge1$. Then
\[
\E H^4\le\exp[-(15/4+64(j-2)/5)d],\quad
I_q(T^2H)^3\le\exp[-(1223/3420+128(j-2)/5)d-5d^2].
\]
The same statements hold with genuine equal-pair factors at the terminal and unit internal root. The following common calculus, including the ordinary fourth-moment transfer, is also used in Section~\ref{proof:R2-C031}.

\paragraph{R29.2 Actual messages and finite calculus}

All operators, inner products, adjoints and expectations use one unchanged
original finite probability measure. Write

\[
Jv=(Ev)\mathbf1,\quad A=J+K,\quad T=J-K=A^*,\quad
B=-K,\quad u=B\mathbf1=-g,
\]
\[
\begin{gathered}
p=1-u,\quad q=1+u,\quad f_n=A^n\mathbf1,\quad\\
\bar f_n=T^n\mathbf1,\quad h_n=Tf_{n-1},\quad W_n=A\bar f_{n-1}.
\end{gathered}
\]

An f arm points entirely outward. An h arm has its first edge toward
the terminal and its remaining edges outward. The actual root unit arm
points outward and supplies p. Every inner edge points from terminal to
root and supplies \(T^L.\) The two outer edges point from root to receiver
and supply \(T^2.\) Thus the bare actual messages are

\[
F=\prod_i h_{a_i}f_{b_i},\qquad H=pT^LF,\qquad U=T^2H.
\tag{R29.2.1}
\]

For R2-C031 the product has exactly one factor, with its stated restrictions.
The root unit arm supplies the factor $p$. Additional root-arm factors are treated in Section~\ref{proof:Q5-C182}.

Put

\[
\begin{gathered}
\delta_i=\|B^i\mathbf1\|_2^2,\quad d=\delta_1,\quad\\
\theta=\frac{13}{32},\quad \kappa=\frac{19}{32},\quad \tau=d-\delta_2.
\end{gathered}
\]

The finite calculus gives

\[
\begin{gathered}
0\le d\le\frac13,\quad d^2\le \delta_2\le\theta d,\quad\\
\|B\|^2\le\theta,\quad \delta_{i+1}\le\theta \delta_i,\quad\\
|Bv(x)|\le\|v\|_2.
\end{gathered}\tag{R29.2.2}
\]

The lower bound \(\delta_2\ge d^2\) also follows from \(EB^2 1=-d\) and Cauchy.
The spectral measure of -\(B^2\) at 1 is a probability measure, so
\(\delta_3\ge d^3.\) 

After zero-mass atoms are removed, antisymmetry gives K(x, \(x)=0,\allowbreak \) whence
p(x), \(q(x)\ge \mu(x)>0.\) Every actual rooted tree message is positive:
assigning all integrated vertices the root's label produces a positive
term with every diagonal factor equal to one. The finite path calculus
also gives \(0<f_n,\allowbreak  \bar f_n,\allowbreak  h_n\le 2\) and \(Ef_n=E\bar f_n\le 1.\)

The spectral measure nu of -\(B^2\) at u has mass d and
\(\int x^r\,\allowbreak \allowbreak  d\nu=\delta_{r+1}\), supported on [0, theta]. At \(z=i\) sqrt(x),
if a real polynomial Q has coordinates \(Q(z)=E(x)+zD(x),\allowbreak \) then

\[
\begin{gathered}
E[Q(B)u]=-\int D(x)\,d\nu(x),\qquad\\
\|Q(B)u\|^2=\int(E(x)^2+xD(x)^2)\,d\nu(x).
\end{gathered}\tag{R29.2.3}
\]


\paragraph{R29.3 All finite path and unequal-pair estimates}

\paragraph{R29.3.1 Finite convolution and a uniform mean loss}\label{part:R29.3.1}

Define \(\beta_1=\beta_2=0\) and

\[
\beta_n=\sum_{r=1}^{\lfloor(n-1)/2\rfloor}(-1)^{r-1}\delta_r
\quad(n\ge3).
\]

The exact finite convolution is

\[
\begin{gathered}
S_n=\sum_{r=0}^nK^r1\\
=f_n+\sum_{r=3}^n\beta_r f_{n-r},\qquad\\
\tau\le\beta_r\le d.
\end{gathered}\tag{R29.3.1}
\]

The same identity holds after full reversal for \(\bar f_n\) and B. The beta coefficients are nonnegative by decreasing alternating moments.

Writing \(\mu_n=Ef_n,\allowbreak \) the exact initial means are

\[
\begin{aligned}
\mu_2&=1-d,&\mu_3&=1-2d,\\
\mu_4&=1-3d+\delta_2,&
\mu_5&=1-4d+2\delta_2+d^2,\\
\mu_6&=1-5d+3\delta_2+3d^2-\delta_3,&
\mu_7&=1-6d+4\delta_2+6d^2-2\delta_3-2d\delta_2.
\end{aligned}
\]

Their losses divided by d, for \(n=4,\allowbreak \) 5, 6, 7, are at least respectively
\(83/32,\allowbreak  137/48,\allowbreak  89/32,\allowbreak  127/48.\) For \(n=7,\allowbreak \) for example, discard the
positive \(2\delta_3\) and use

\[
6d-(4-2d)\delta_2-6d^2
\ge[6-4\theta-(6-2\theta)d]d
\ge\frac{127}{48}d.
\]

For every \(n\ge 8,\allowbreak \) retain the six actual terms corresponding to
\(\mu_0,\allowbreak ...,\allowbreak  \mu_5\) in (R29.3.1). The three coefficients \(\beta_{n+1},\allowbreak  \beta_n,\allowbreak  \beta_{n-1}\) are at least tau+kappa \(\delta_3.\) Consequently

\[
\begin{aligned}
1-\mu_n
&\ge\tau(7-10d+3\delta_2+d^2)+3\kappa \delta_3\\
&\ge\kappa d(7-10d+7d^2)
\ge\frac{95}{36}d>\frac{83}{32}d.
\end{aligned}
\tag{R29.3.2}
\]

This proves \(\mu_n\le 1-83d/32\) for every \(n\ge 4.\)

\paragraph{R29.3.2 The difference formula and every unequal pair}

For \(a<b\) let \(R_{ab}=\sum_{r=a+1}^bK^r1\). Subtracting two finite
convolutions gives exactly

\[
\begin{gathered}
f_b-f_a=R_{ab}\\
+\sum_{i=0}^{a-3}(\beta_{a-i}-\beta_{b-i})f_i\\
-\sum_{i=\max(0,a-2)}^{b-3}\beta_{b-i}f_i.
\end{gathered}\tag{R29.3.3}
\]

Empty sums are zero. At $a=1$ the identity reads

\[
f_b-f_1=R_{1b}-\sum_{i=0}^{b-3}\beta_{b-i}f_i.
\tag{R29.3.3a}
\]

Since \(h_1+f_1=2,\allowbreak \) retaining the genuine \(f_0=1\) term gives
\(h_1+f_b\le 2+R_{1b}-\beta_b\) for \(b\ge 3.\) For \(b=2\) the sum is empty
and \(\beta_2=0,\allowbreak \) yielding exactly the same bound. AM--GM and
log \(t\le t-1\) then give \(h_1f_b\le \exp(R_{1b}-\beta_b).\)
This proves the $a=1$ row.


In the first sum of (R29.3.3) a positive
coefficient requires an odd starting alternating-moment index; its first
omitted even-index moment is \(\delta_{2r}.\) Each such first moment occurs at
most twice, and \(f_i\le 2.\) Therefore the positive part of that entire sum
is at most \(4\delta_2/(1-\theta^2)\). For \(a\ge 5,\allowbreak \) (R29.3.2) and
\(f_a+h_a=2\mu_{a-1}\) yield

\[
h_a+f_b\le2+R_{ab}
-\left(\frac{83}{16}-\frac{1664}{855}\right)d,
\]
\[
\frac{83}{16}-\frac{1664}{855}
=\frac{16}{5}+\frac{113}{2736}.
\tag{R29.3.4}
\]

For \(a=4\) retain the actual negative term -tau \(f_2\) in (R29.3.3):

\[
f_b-f_4\le R_{4b}+\delta_2(1+p)-\tau f_2.
\]

The original row identity

\[
2Kg(x)=E_{y,z}[(K(x,y)-K(x,z))K(y,z)]
\]

and |\(a-b|\le 1-ab\) for a, b in [-1, 1] give |\(Kg|\le pq/2.\)
Hence \(f_2=p+Kg\ge p^2/2,\allowbreak \) and completing this real quadratic gives

\[
\begin{gathered}
\delta_2(1+p)-\tau f_2\\
\le \delta_2+\frac{\delta_2^2}{2\tau}\\
\le\frac{663}{1216}d<\frac34d.
\end{gathered}\tag{R29.3.5}
\]

The \(a=4\) pair sum therefore loses at least 13d/4.
For any positive pair X, Y, \(XY\le((X+Y)/2)^2\) and
\(\log t\le t-1\) give \(\log(XY)\le X+Y-2\).
Equations (R29.3.4)-(R29.3.5) establish the common strong pair bound

\[
h_af_b\le\exp\left(R_{ab}-\frac{16}{5}d\right)
\qquad(4\le a<b).
\tag{R29.3.6}
\]

The short finite differences give, with no additional arm,

\[
\begin{gathered}
h_1f_b\le e^{R_{1b}-\beta_b},\quad\\
h_2f_b\le e^{R_{2b}-\beta_b},\quad\\
h_3f_b\le e^{R_{3b}-(2d-\delta_2)}.
\end{gathered}\tag{R29.3.7}
\]

For \(a=3,\allowbreak  b\ge 4,\allowbreak \) retain -\(\beta_{b-1}f_1\le -\tau\) p:

\[
h_3f_b\le
\exp\{R_{3b}+\tau u-(3d-2\delta_2)\}.
\tag{R29.3.8}
\]

For the exact pair (2, 3),

\[
h_2+f_3=2-d+K^31,\qquad h_2f_3\le e^{-d+K^31}.
\tag{R29.3.9}
\]

The retained tau u in (R29.3.8) is the scalar correction for the same kernel.

\paragraph{R29.4 The signed-mean exponential certificate}\label{calc:five-signed-exp}
The full signed-mean lemma of Section~\ref{proof:Q5-C081},
paragraph~\ref{part:Q81.2}, gives
\begin{equation}\tag{R29.4.1}\label{eq:oc:r29-signed}
 \E e^{4BV}\le e^{4\E BV+(17/4)\|V\|^2}.
\end{equation}
Its scalar statement is
\begin{equation}\tag{R29.4.2}\label{eq:oc:r29-scalar-signed}
 \E e^{sW}\le e^{s\E W+(17/64)s^2U^2}
 \quad(|W|\le U,\ \E W^2\le\theta U^2,\ s\ge0).
\end{equation}
These are the signed-mean inequalities proved in
paragraph~\ref{part:Q81.2}. The complete-vector payment is
\eqref{eq:unit-two-step:moment-payment}.

Use Lemmas~\ref{lem:oc:finite-continuation} and~\ref{lem:oc:fourth}.

\paragraph{R29.6 R2-C029: exhaustive finite masks and all later pairs}\label{part:R29.6}

\paragraph{R29.6.1 One mask for every finite number of pairs}

Put

\[
C_\beta=\frac1{1-\theta^2}=\frac{1024}{855},\qquad
C_*=\frac{32768}{16245}.
\]

For every finite polynomial P with coefficients in \{0, 1\}, separate
the positive and negative terms of each real/imaginary coordinate at
i sqrt(x). Each coordinate's absolute value is bounded by the larger
one-sided geometric tail, not the sum of both tails. This gives

\[
\begin{gathered}
|P(i\sqrt x)|^2\le\frac{1+x}{(1-x^2)^2}\\
\le C_*<\frac{101}{50}<\frac{100}{49}.
\end{gathered}\tag{R29.6.1}
\]

The exponent intervals [\(a_i+1,\allowbreak  b_i\)] are disjoint because the actual
arm lengths are strictly ordered. Their union therefore makes one
finite 0-1 mask, regardless of j or any gap. The geometric tails bound these finite sums with a norm bound independent
of j.

Useful exact constants are

\[
\begin{gathered}
C_\beta\theta=\frac{416}{855}<\frac12,\quad\\
C_\beta\theta^2=\frac{169}{855}<\frac15,\quad\\
C_\beta\theta^3=\frac{2197}{27360}<\frac1{12},\quad\\
C_\beta\theta^4=\frac{28561}{875520}<\frac1{30}.
\end{gathered}\tag{R29.6.2}
\]

Also \(\theta^2<1/6,\allowbreak  \theta^3<1/14,\allowbreak \) and \(\theta^{3/2}<4/15.\)
Let \(\Delta=(16/5)(j-2)d\). The following five cases are a
structural partition of every \(j\ge 2\) table:

\begin{table}[htbp]
\centering\small
\setlength{\tabcolsep}{5pt}
\caption{The five finite even-terminal packets, with $R_0=\sum_iR_{a_i b_i}$ and $\Delta=(16/5)(j-2)d$.}\label{tab:oc:r29-structure}
\begin{tabular}{clcc}
\toprule
Packet & Arm condition & $R$ & Lower $\Gamma$\\
\midrule
E2 & $a_1=2$ & $R_0$ & $(607/160)d+\Delta$\\
E3 & $a_1\ge3$ & $R_0$ & $(767/160)d+\Delta$\\
I & $a_1=1,\ b_1\ge3$ & $R_0$ & $(607/160)d+\Delta$\\
II & $a_1=1,\ b_1=2,\ a_2\ge4$ & $R_0$ & $(16/5)d+\Delta$\\
III & $a_1=1,\ b_1=2,\ a_2=3$ & $R_0+\tau u$ & $3d-2\delta_2+\Delta$\\
\bottomrule
\end{tabular}
\end{table}

For E2 and I the first pair loses at least kappa d and the next at
least 16d/5. For E3 the first loses at least \(2d-\delta_2\) and the next at
least 16d/5; \(a_1\ge 4\) only improves that. Case II discards the
nonnegative first-pair loss and keeps the next strong pair. Case III
uses (R29.3.8) on its actual \(a_2=3\) pair. Every pair after the second has
shorter length at least four and supplies exactly the extra Delta
loss used here. This justifies the linear increment in j without
discarding any later arm.

\paragraph{R29.6.2 Complete scalar packets}

The following table records proved bounds, with entries normalized by d.
The row \(D_2\) for III retains its actual sharper value \(\delta_2.\)

\begin{table}[htbp]
\centering\small
\setlength{\tabcolsep}{5pt}
\caption{Scalar packet bounds for Lemma~\ref{lem:oc:fourth}. Interval rows contain $\delta/d$ and $\epsilon/d$; other rows are upper bounds. The $D_2$ entry in III retains $\delta_2$, a moment of the same host.}\label{tab:oc:r29-packets}
\begin{tabular}{lccccc}
\toprule
Quantity & E2 & E3 & I & II & III\\
\midrule
$c$ & $5/4$ & $6/5$ & $3/2$ & $3/2$ & $4/3$\\
$\sigma/d$ & $1/3$ & $9/64$ & $325/392$ & $325/512$ & $1277/2592$\\
$\delta/d$ & $[-1/5,91/100]$ & $[-1/5,2/3]$ & $[-6/5,3/4]$ & $[-6/5,1/25]$ & $[-1,3/20]$\\
$\epsilon/d$ & $[-7/10,7/10]$ & $[-7/8,7/8]$ & $[-17/10,3/4]$ & $[-13/10,6/25]$ & $[-4/3,7/20]$\\
$D_1/d$ & $0$ & $1/5$ & $0$ & $1/5$ & $1/5$\\
$D_2$ & $d/12$ & $d/12$ & $d/2$ & $d/2$ & $\delta_2$\\
$D_3/d$ & $1/5$ & $1/30$ & $1/5$ & $1/5$ & $1/8$\\
$D_k/d,\ k\ge4$ & $1/12$ & $1/12$ & $1/5$ & $1/5$ & $1/12$\\
$h$ & $1/2$ & $1/2$ & $3/8$ & $3/8$ & $1/2$\\
\bottomrule
\end{tabular}
\end{table}

Here is a reconstruction of each packet, including the infinite index
quantifier on \(D_k.\)

E2. Write \(R=-B^2\) P(-B)u, where \(P_0=1,\allowbreak \) and
\(P(-z)=U-zV.\) The coordinates of R are xU,-xV.
One-sided alternating tails give

\[
\frac12<U<\frac65,\qquad -C_\beta x\le V\le C_\beta.
\]

Therefore

\[
D_0=\int xV\in[-d/5,d/2],\qquad
D_1=-\int xU\in[-d/2,0].
\]

The row predecessor and the spectral norm give

\[
\begin{gathered}
|R|^2\le C_*\theta d<\frac9{25},\qquad\\
\sigma\le C_*\theta^2d=\frac{5408}{16245}d<\frac d3.
\end{gathered}
\]

Thus \(c=5/4\) is valid. The positive delta bound is

\[
\begin{gathered}
\delta\le\\
\left(\frac{416}{855}+\frac54\frac{5408}{16245}\right)d\\
=\frac{4888}{5415}d<\frac{91}{100}d.
\end{gathered}
\]

The useful epsilon bound keeps \(D_0+D_1\) together. Since
\(V-U\le x/(1-x),\allowbreak \) one has

\[
\begin{gathered}
D_0+D_1\le\frac{13}{19}\delta_2,\\
\qquad \epsilon\le\frac{13}{19}\delta_2+\frac54\sigma<\frac7{10}d.
\end{gathered}
\]

The negative bound is \(D_0+D_1\ge -7d/10.\) Multiplying the two
coordinates successively by z shows the positive part of \(D_2\) begins
at \(\delta_4,\allowbreak \) that of \(D_3\) at \(\delta_3,\allowbreak \) and those of every \(k\ge 4\) no earlier than
\(\delta_4.\) Within each of the four power classes, subsequent positive
even-moment indices increase by two, so (R29.6.2) bounds the entire finite
tail. In particular \(\gamma_3\ge 11d/12,\allowbreak  \gamma_4\ge 4d/5,\allowbreak \) and every
\(\gamma_k\) for \(k\ge 5\) is at least (\(\kappa-1/12)d=49d/96.\) Hence \(h=1/2.\)

E3. Write \(R=B^3P(-B)u,\allowbreak \) allowing \(P_0=0.\) The general coordinates
U, V lie between -\(C_\beta\) x and \(C_\beta\); those of R are
-\(x^2V,\allowbreak -xU.\) Consequently

\[
\begin{gathered}
\sigma\le C_*\theta^3d=\frac{2197}{16245}d<\frac9{64}d,\\
\quad |R|^2\le C_*\theta^2d<\frac13d\le\frac19.
\end{gathered}
\]

Use \(c=6/5.\) The signed means are

\[
 D_0\in[-C_\beta \delta_3,C_\beta \delta_2],\qquad
 D_1\in[-C_\beta \delta_4,C_\beta \delta_3].
\]

Thus \(\delta\ge -d/5\) and
\(\delta\le 52702d/81225<2d/3.\) Epsilon lies between -17d/60 and
13d/15, which is inside the stated [-\(7/8,\allowbreak  7/8\)] interval.
The positive \(D_2\) tail begins at \(\delta_4,\allowbreak  D_3\) at \(\delta_5,\allowbreak \) and every later
positive tail no earlier than \(\delta_4.\) Four-power continuation and (R29.6.2)
give the table. In particular \(\gamma_3\ge 11d/12,\allowbreak  \gamma_4\ge 29d/30,\allowbreak \) and the remaining \(\gamma_k\ge 49d/96\); \(h=1/2.\)

I. Here \(R=BP(-B)u\) and \(P_0=P_1=1.\) Its coordinates are xV, U,
where

\[
1-C_\beta x\le U,V\le C_\beta,\qquad \frac12<U,V<\frac65.
\]

The row bound gives |\(R|<1\); (R29.6.1) gives
\(\sigma\le (325/392)d.\) Use \(c=3/2.\) The actual first two moments are

\[
D_0=-\int U\in[-6d/5,-439d/855],\qquad
D_1=-\int xV\in[-d/2,0].
\]

Hence \(\delta\le (-439/855+(3/2)(325/392))d<3d/4,\allowbreak \)
and the other delta/epsilon bounds follow without replacing signs.
\(D_2\le C_\beta \delta_2<d/2,\allowbreak  D_3\le C_\beta \delta_3<d/5,\allowbreak \) and each \(k\ge 4\)
has positive part bounded by d/5 through the same four-power
continuation. Thus \(\gamma_3\ge d/2,\allowbreak  \gamma_4\ge 4d/5,\allowbreak \) and all later
\(\gamma_k\ge (\kappa-1/5)d=63d/160>3d/8.\)

II. The actual missing exponents impose \(P_0=1,\allowbreak  P_1=P_2=0.\)
Consequently

\[
\begin{gathered}
1-C_\beta x^3\le U\le C_\beta,\qquad\\
-C_\beta x\le V\le C_\beta x^2,\\
\quad U\ge\frac{11}{12}.
\end{gathered}
\]

The sharper mask bound is

\[
\begin{gathered}
|P|^2\le\frac{1+x^3}{(1-x^2)^2}\\
\le\frac{24864}{16245}<\frac{54}{35}<\frac{25}{16}.
\end{gathered}
\]

This yields \(\sigma\le (325/512)d\) and |\(R|<1\); take \(c=3/2. D_0\) is in [-6d/5,-11d/12] and
\(D_1\) is in [-d/12, d/5]. Therefore

\[
\delta\le-\frac{11}{12}d+\frac{975}{1024}d
=\frac{109}{3072}d<\frac1{25}d.
\]

The lower epsilon bound is -13d/10 and its upper bound is 6d/25.
Also \(D_2\le \delta_2+C_\beta \delta_4<d/2.\) The remaining means have
the same sufficient upper bounds as I. The full gamma tail has
the same \(h=3/8.\) Missing original powers are used only as mask
information; no shorter path replaces the finite original path.

III. Now \(P_0=P_2=1\) and \(P_1=0.\) Its finite coordinates satisfy

\[
\begin{gathered}
1-C_\beta x\le U\le1-\frac x3\le1,\quad\\
-C_\beta x\le V\le C_\beta x^2,\quad |V|\le\frac43x.
\end{gathered}
\]

Indeed \(U\le 1-x+C_\beta x^2\le 1-x/3\) on the full range. Therefore

\[
\begin{gathered}
|P|^2\le(1-x/3)^2+\frac{16}{9}x^3\le1-\frac x6\le1\\
\quad(0\le x\le1/2).
\end{gathered}\tag{R29.6.3}
\]

Here \(R=R_0+\tau\) u. The complete row predecessor is P(-B)u+tau 1.
Its constant-vector cross term is retained:

\[
\begin{gathered}
|R|^2\le d+\tau^2+2\tau E[P(-B)u]\\
\le d+\tau^2+\frac{2\tau}{5}d\\
\le\frac{22}{45}<\frac9{16}.
\end{gathered}\tag{R29.6.4}
\]

The bound on the cross mean follows from
E[P(-B)u]=int \(V\le C_\beta \delta_3<d/5.\) Thus the globally valid
Taylor choice is \(c=4/3.\) In the squared norm of R itself,
\(E(uR_0)=int xV\le C_\beta \delta_4<d/12,\allowbreak \) so

\[
\begin{gathered}
\sigma\le \delta_2+\frac{\tau}{6}d+\tau^2d\\
\le\left(\frac{13}{32}+\frac1{27}+\frac4{81}\right)d\\
=\frac{1277}{2592}d.
\end{gathered}\tag{R29.6.5}
\]

Since \(D_0=-int\) U lies in [-d,-439d/855],

\[
-d\le\delta\le
\left(-\frac{439}{855}+\frac{1277}{1944}\right)d
<\frac3{20}d.
\]

\(D_1=-E(uR_0)-\tau\) d lies above -d/12-tau d and below d/5.
The resulting epsilon lower bound is -\(47d/36>-4d/3\) and
the upper bound is 7d/20. The retained finite moments give

\[
\begin{gathered}
D_2\le \delta_2-\delta_3+C_\beta \delta_4\le \delta_2,\quad\\
D_3\le C_\beta \delta_5+\tau \delta_2\\
<\left(\frac1{30}+\frac{13}{144}\right)d\\
=\frac{89}{720}d<\frac18d.
\end{gathered}
\]

The all-k tail is essential. For every even \(k\ge 4,\allowbreak \) tau u has zero
mean after \(B^k\) and the positive \(R_0\) tail starts at \(\delta_4.\) For \(k=5\) mod 4,
the tau term is nonpositive and the positive \(R_0\) tail starts at \(\delta_5.\)
For \(k=7\) mod 4, the first possible positive contributions are
\(C_\beta \delta_7\) and tau \(\delta_4\); together they are less than
(\(1/12+2/9)\delta_4<11d/504<d/30.\) Moving four powers onward shifts
these even-moment indices by two and only reduces the tail bound.
Thus \(D_k\le d/12\) for every \(k\ge 4.\) Accordingly
\(\gamma_3\ge \kappa\) d, \(\gamma_4\ge 7d/8\) and all later \(\gamma_k\ge 49d/96\);
\(h=1/2\) is valid for every integer index.

This establishes the complete packet table. In particular, the dangerous
III row retains both its actual constant-vector cross and its stronger
long-tail gamma value.

\paragraph{R29.7 R2-C029: all four inner-length ranges}

In this section every displayed lower loss coefficient omits the
additional 4Delta, which is restored at the end.

\paragraph{R29.7.1 \(L=1\)}

Substitute the packets into (R29.5.6). Five valid loss coefficients are

\begin{table}[htbp]
\centering\small
\setlength{\tabcolsep}{5pt}
\caption{Lower loss coefficients for $L=1$.}\label{tab:oc:r29-L1}
\begin{tabular}{cl}
\toprule
Packet & Lower $-\log\E H^4/d$\\
\midrule
E2 & $607/40-14/5-17/12$\\
E3 & $767/40-7/2-153/256$\\
I & $607/40-3-119/32$\\
II & $64/5-24/25-17/6$\\
III & $35/4-7/5-17/8=209/40$\\
\bottomrule
\end{tabular}
\end{table}

Each exceeds \(15/4.\) The weaker sigma bounds 7d/8, 2d/3, d/2
used in the last three rows follow from the already proved packets.
All signed epsilon terms have their proper upper bounds.

\paragraph{R29.7.2 \(L=2\)}

Let \(V_2=u+BR.\) For E2,
\(\lVert BR\rVert ^2\le 2197d/16245<(37/100)^2d,\allowbreak \) so
\(\lVert V_2\rVert \le 11\sqrt{d}/8.\) For E3, \(\lVert BR\rVert \le \sqrt{d}/4,\allowbreak \)
so \(\lVert V_2\rVert \le 5\sqrt{d}/4.\)

For I and II its coordinates are 1-xU, xV. Since \(U>1/2\) and
|\(P|^2<100/49,\allowbreak \)

\[
|1-xP(-z)|^2\le1-x+\frac{100}{49}x^2\le1
\quad(0\le x\le\theta).
\]

For III its coordinates are 1-xU, tau+xV. The first lies between
zero and 1-x/2; the absolute value of the second is below \(1/3,\allowbreak \)
using \(\tau\le 2/9,\allowbreak \) -\(C_\beta x^2\le xV\le C_\beta x^3.\)
Hence its squared modulus is at most
(\(1-x/2)^2+x/9\le 1.\) These estimates keep tau as its actual scalar.

Apply (R29.5.7), using respectively |\(\delta|\le d,\allowbreak \) 2d/3, 6d/5, 6d/5, d:

\begin{table}[htbp]
\centering\small
\setlength{\tabcolsep}{5pt}
\caption{Lower loss coefficients for $L=2$.}\label{tab:oc:r29-L2}
\begin{tabular}{cl}
\toprule
Packet & Lower $-\log\E H^4/d$\\
\midrule
E2 & $607/40-14/5+4-1/3-2057/256-10/9$\\
E3 & $767/40-7/2+4-1/3-425/64-40/81$\\
I & $607/40-3+4-2-17/4-8/5$\\
II & $64/5-24/25+4-2-17/4-8/5$\\
III & $35/4-7/5+4-13/8-17/4-10/9=1571/360$\\
\bottomrule
\end{tabular}
\end{table}

All exceed \(15/4\); III alone has the explicit margin \(221/360.\)
The -4d mean and \(D_2\) terms in (R29.5.7) are retained in every row.

\paragraph{R29.7.3 \(L=3,\allowbreak \) including the continuous III comparison}

Write \(V_3=(1+B)u+\delta u+B^2R.\) The mean part
\(4\epsilon+4D_2+4D_3-4\delta\) d has the following upper coefficients,
and the curvature term (\(10/3)\epsilon^2\) has the listed upper
coefficients after using \(d^2\le d/3:\)

\begin{table}[htbp]
\centering\small
\setlength{\tabcolsep}{5pt}
\caption{Upper signed-mean and curvature costs at $L=3$.}\label{tab:oc:r29-L3-costs}
\begin{tabular}{ccc}
\toprule
Packet & Mean cost divided by $d$ & Curvature cost divided by $d$\\
\midrule
E2 & $21/5$ & $49/90$\\
E3 & $127/30$ & $245/288$\\
I & $37/5$ & $289/90$\\
II & $134/25$ & $169/90$\\
III & $583/120$ & $160/81$\\
\bottomrule
\end{tabular}
\end{table}

The negative lower bounds on delta are needed for the term -4delta d.
They have not been dropped.

For E2, triangle bounds give

\[
\frac{\|V_3\|}{\sqrt d}
<\frac{19}{16}+\frac{91}{300}+\frac14<\frac74.
\]

For E3 the corresponding sum is
\(19/16+2/9+1/6<8/5.\) The exceptional \(G_2\) norm is explicitly
\(\sqrt{1+\theta}<19/16\); it is not bounded by one.

For I and II the coordinates are

\[
A_3=1+\delta-x^2V,\qquad C_3=1-xU.
\]

In I, both are positive, \(A_3\le 5/4\) and \(C_3\le 1,\allowbreak \) so
\(\lVert V_3\rVert ^2\le (63/32)d.\) In II,
\(0<A_3\le 329/300\) and \(0<C_3\le 1-x/2.\) Since
\(x(1-x/2)^2\le 9/32\) on [0, \(1/2\)],

\[
\|V_3\|^2\le[(329/300)^2+9/32]d<\frac32d.
\]

For III the exact coordinates are

\[
A_3=1+\delta-\tau x-x^2V,\qquad C_3=1-xU.
\tag{R29.7.1}
\]

The lower bounds give \(A_3>1/2\) and \(C_3\ge \kappa.\) Their upper bounds,
with the same original delta and tau, are

\[
A_3\le1+d(3/20-\kappa x)+x^2/2,\qquad C_3\le1-x/2.
\tag{R29.7.2}
\]

For \(0\le x\le 1/4,\allowbreak \) (R29.7.2) gives

\[
A_3^2+xC_3^2
\le(173/160)^2+\frac14
=\frac{36329}{25600}<\frac{23}{16},
\]

with difference \(471/25600.\) For \(1/4\le x\le \theta,\allowbreak  3/20-\kappa x\le 1/640,\allowbreak \) so set \(\epsilon_0=1/1920\) and obtain

\[
A_3^2+xC_3^2
\le1+x+x^3/4+x^4/4+
2\epsilon_0+\epsilon_0x^2+\epsilon_0^2.
\]

The last three terms are less than \(1/800.\) Thus the entire expression
is less than
\(1+13/32+1/56+1/128+1/800,\allowbreak \) which is below \(23/16\) by \(97/22400.\)
This is a proof on the entire continuous spectral interval, not
checking selected frequencies or assuming \(x\ge d.\) It gives
\(\lVert V_3\rVert ^2\le (23/16)d.\)

Substitution into (R29.5.8) gives

\begin{table}[htbp]
\centering\small
\setlength{\tabcolsep}{5pt}
\caption{Lower loss coefficients for $L=3$.}\label{tab:oc:r29-L3}
\begin{tabular}{cl}
\toprule
Packet & Lower $-\log\E H^4/d$\\
\midrule
E2 & $607/40+8-21/5-49/90-833/64>5$\\
E3 & $767/40+8-127/30-245/288-272/25>11$\\
I & $607/40+8-37/5-289/90-1071/128=24173/5760$\\
II & $64/5+8-134/25-169/90-51/8>7$\\
III & $35/4+8-583/120-160/81-391/64=98677/25920$\\
\bottomrule
\end{tabular}
\end{table}

The I margin over \(15/4\) is \(2573/5760\); the III margin is
\(1477/25920.\) These are positive exact rational margins.

\paragraph{R29.7.4 Every finite \(L\ge 4\)}

For every finite \(m\ge 3,\allowbreak \) \(\lVert G_m(B)\rVert \le 1.\) To see the complete length
scope, at \(z=i\) sqrt(x), an odd m gives
|\(G_m|^2=(1+x^m)/(1+x)\le 1.\) For m divisible by four the numerator is
(\(1-x^{m/2})^2.\) For \(m=2\) mod 4 and \(m\ge 6\) it is at most
(\(1+x^3)^2\le 1+x\) for \(x\le 1/2.\) The omitted \(m=2\) is the genuine
exception already used correctly in \(L=3.\)

Every packet has \(\sigma\le 7d/8.\) Hence

\[
\begin{gathered}
\|B^{L-1}R\|^2\le\theta^3\sigma<d/16,\qquad\\
\|V_L\|\le\frac54\sqrt d.
\end{gathered}\tag{R29.7.3}
\]

The gamma hypotheses were proved for every index in \mbox{\hyperref[part:R29.6]{part R29.6}},
not merely for \(k\le L.\) Equation (R29.5.10) now has norm cost 425d/64
and curvature cost at most \(10h^2d/9.\) Its five lower loss coefficients
are

\begin{table}[htbp]
\centering\small
\setlength{\tabcolsep}{5pt}
\caption{Lower loss coefficients for every $L\ge4$.}\label{tab:oc:r29-long}
\begin{tabular}{cl}
\toprule
Packet & Lower $-\log\E H^4/d$\\
\midrule
E2 & $607/40+6-182/25-425/64-5/18$\\
E3 & $767/40+6-16/3-8/5-425/64-5/18$\\
I & $607/40+9/2-6-425/64-5/32$\\
II & $64/5+9/2-8/25-8/5-425/64-5/32$\\
III & $35/4+6-6/5-8/5-425/64-5/18$\\
\bottomrule
\end{tabular}
\end{table}

Each is greater than 5, and therefore greater than \(15/4.\) These
are one finite-polynomial proof for all \(L\ge 4,\allowbreak \) with no cutoff or
unproved limiting passage.

Restoring the omitted 4Delta in all four length ranges proves

\[
\begin{gathered}
E[(pT^LF)^4]\le\\
\exp\left[-\left(\frac{15}{4}+\frac{64}{5}(j-2)\right)d\right].
\end{gathered}\tag{R29.7.4}
\]

This completes R2-C029's ordinary-fourth claim for all five structural
classes, every \(j\ge 2\) and every original \(L\ge 1.\)


\paragraph{R29.11 The exact C094 ordinary-fourth bridge}\label{part:R29.11}

The input to the Q5-C094 Theorem A, Sections 2.1--2.2, transfer is every original real $Y$ with its ordinary fourth moment $E|Y|^4$. Its same-kernel row and norm argument is reproduced here.

\paragraph{R29.11.1 Same-kernel inverse fourth of the genuine \(\bar f_2\)}

For this subsection use the original \(g=K1\) and \(r=Kg=K^2\) 1.
Then \(\bar f_2=q+r,\allowbreak  Er=-d,\allowbreak  Egr=0,\allowbreak  Er^2=\delta_2.\) The original centered-row
Cauchy inequality, valid for every mean-zero v, is

\[
|Kv(x)|^2\le p(x)q(x)E[v^2].
\tag{R29.11.1}
\]

Indeed subtract g(x) from the row and use
E[(K(x, \(y)-g(x))^2\)]\ensuremath{\le}\(1-g(x)^2.\) Applying this to \(v=g\) gives
\(r^2\le pqd.\) This controls the actual small-q denominator:

\[
E[r^4/q]\le dE[pr^2]\le2d\delta_2.
\tag{R29.11.2}
\]

Let \(W=2r-g=K(2g-1)\) and \(t_0=1+2d.\) The original row bound gives
|\(W|\le \sqrt{1+4d}\le t_0,\allowbreak  EW=-2d\) and \(EW^2=d+4\delta_2.\) The pointwise
polynomial (\(W+2d)^2(W-t_0)\le 0\) implies

\[
EW^3\le(t_0-4d)(d+4\delta_2)-4t_0d^2+8d^3.
\]

The exact cubic combination in \(\bar f_2^4/q\) is

\[
\mathcal C=-Eg^3+4E(g^2r)-6E(gr^2)+4Er^3
=\frac{EW^3+E(g^2W)}2.
\]

Using \(E(g^2W)\le t_0d\) gives

\[
\mathcal C\le d-2d^2+2(1-2d)\delta_2.
\]

The full expansion, with \(Egr=0\) and the same \(\delta_2\) retained, is

\[
E[\bar f_2^4/q]
=1-d+6\delta_2+\mathcal C+E[r^4/q]
\le1-2d^2+(8-2d)\delta_2.
\tag{R29.11.3}\label{eq:unit-output:inverse-row}
\]

There is no independent choice of moment extrema in this expansion
and no host-independent lower bound on q.

\paragraph{R29.11.2 Exact transfer exponents and the original norm}

For any positive original kernel operator \(\mathcal B\), let
\(h= \mathcal B1\). Original-row Cauchy and then Hölder with
exponents 3 and \(3/2\) give

\[
|\mathcal BY|^2\le h\mathcal B(|Y|^2),
\]
\[
\begin{aligned}
I_q(\mathcal BY)
&\le E[(h^4/q)^{1/3}
       (\mathcal B(|Y|^2)^2)^{2/3}]\\
&\le E[h^4/q]^{1/3}
       E[\mathcal B(|Y|^2)^2]^{2/3}.
\end{aligned}
\]

Hence, in the unchanged ordinary L2(mu) space,

\[
I_q(\mathcal BY)^3
\le E[h^4/q]\|\mathcal B\|^4(E|Y|^4)^2.
\tag{R29.11.4}\label{eq:unit-output:ordinary-transfer}
\]

For \(\mathcal B=T^2\), the actual h is \(\bar f_2,\allowbreak \) not q.
The exact norm input recorded in C094 is
\(\|T^2\|\le(1+(1024/855)\delta_2)^2\).
This is an accepted original ordinary norm bound; neither its
adjoint nor its input is changed by a q-tilt. Together with (R29.11.3),

\[
I_q(T^2Y)^3\le B_2(d,\delta_2)(E|Y|^4)^2,
\]
\[
B_2=\{1-2d^2+(8-2d)\delta_2\}
\left(1+\frac{1024}{855}\delta_2\right)^8.
\tag{R29.11.5}
\]

This transfer uses the stated original norm estimate for $T^2$.

The negative curvature is also retained. Since
\(8-2d>0,\allowbreak \) the first factor is at most

\[
R=1+\frac{13}{4}d-\frac{45}{16}d^2.
\]

If \(v=R-1,\allowbreak \) then \(0\le v\le 37/48\) and \(v\ge 37d/16.\) Integrating
1/(1+t), or differentiating the difference, gives
\(\log(1+v)\le v-v^2/\)[2(1+v)]. It follows that

\[
\log R\le\frac{13}{4}d-\frac{11757}{2720}d^2.
\]

Let \(\lambda_0=(1024/855)\theta=416/855,\allowbreak \) with
\(12/25<\lambda_0<1/2.\) Applying the same logarithm inequality,

\[
\begin{gathered}
8\log(1+\lambda_0d)\\
\le8\lambda_0d-\frac{4\lambda_0^2}{1+\lambda_0/3}d^2\\
\le8\lambda_0d-\frac34d^2.
\end{gathered}
\]

Here (\(24/7)(12/25)^2=3456/4375>3/4,\allowbreak \) while
\(11757/2720+3/4=13797/2720>5.\) Therefore

\[
B_2(d,\delta_2)\le
\exp\left(\frac{24427}{3420}d-5d^2\right).
\tag{R29.11.6}\label{eq:unit-output:curvature}
\]

For R2-C029 use (R29.7.4); for R2-C031 use the ordinary fourth bound in Section~\ref{proof:R2-C031}. Since
\(15/2-24427/3420=1223/3420,\allowbreak \) they give respectively

\[
I_q(T^2H)^3\le
\exp\left[-\left(\frac{1223}{3420}+
\frac{128}{5}(j-2)\right)d-5d^2\right],\tag{R29.11.7}
\]
\[
I_q(T^2H)^3\le
\exp\left[-\frac{1223}{3420}d-5d^2\right].
\tag{R29.11.8}
\]

The transfer uses the ordinary fourth moment of $H$ in the original $L^2(\mu)$ space.

\paragraph{R29.12 Genuine decorations, full reversal, and boundary cases}

\paragraph{R29.12.1 Every original equal pair remains}

At each original terminal or internal root, pair equal-length arms
as far as possible. Each pair has its two actual directions \(f_l,\allowbreak  h_l,\allowbreak \)
and

\[
0<f_lh_l\le(Ef_{l-1})^2\le1.
\]

Let \(E_t\) and \(E_r\) be the actual products of all these pairs at their
own vertices. The decorated component is exactly

\[
H_D=E_r pT^L(E_tF),\qquad U_D=T^2H_D.
\tag{R29.12.1}
\]

Positivity proves \(H_D\le pT^LF\) and \(U_D\le T^2(pT^LF),\allowbreak \) so both
ordinary-fourth and Iq estimates persist. The terminal pairs remain
inside \(T^L,\allowbreak \) the root pairs remain outside it at the actual root,
and every pair's edges remain in the tree. 

\paragraph{R29.12.2 Complete reversal}

Replacing K by -K reverses every actual edge simultaneously. It
interchanges $A,T$; $p,q$; $f,\bar f$; and $h,W$. Thus the fully reversed
component is

\[
A^2\left[E_r^-qA^L\left(E_t^-\prod_iW_{a_i}\bar f_{b_i}\right)\right],
\tag{R29.12.2}
\]

with a single pair in R2-C031, and the corresponding budget is \(I_p.\)
The decorations \(E_t^-,\allowbreak  E_r^-\) are the actual \(\bar f_lW_l\) products.
The norm moments \(\delta_i\) remain the same under reversal.
This second direction is fixed by the structure before observing K. Reversing only the inner component
would not justify keeping the same Iq claim. Both outer edges in
(R29.12.2) are reversed as well.

\paragraph{R29.12.3 Small masses, extreme kernels and \(d=0\)}

All estimates allow arbitrary finite positive atom masses, including
arbitrarily small ones, and pointwise K values at either endpoint
of [-1, 1]. The positive diagonal contribution guarantees that the
actual messages and p, q are positive on the retained support.
No estimate requires a mass lower bound uniform over hosts.
The only inverse q expression used for the output bridge is paid
by the actual centered-row inequality (R29.11.1)-(R29.11.2).

When \(d=0,\allowbreak \) only \(K1=0\) follows. Then \(p=q=1\) and all actual directed
tree messages equal one by recursively integrating their leaves.
This works for a nonzero skew K with \(K1=0.\) All displayed actual
moment and Iq estimates are equalities at \(d=0.\) The preceding
divisions by d or tau were used only when \(d>0,\allowbreak \) where
\(\tau\ge \kappa d>0.\) 

\paragraph{R29.12.4 Original edge and vertex counts}

Let \(S_pair\) be the sum of one arm length from each genuine equal pair
at the two vertices. Include the receiver c in the component. Then

\[
e_B=L+3+\sum_{i=1}^j(a_i+b_i)+2S_{\rm pair}
\quad\text{for R2-C029},
\tag{R29.12.3}
\]
\[
e_B=L+3+a+b+2S_{\rm pair}
\quad\text{for R2-C031}.
\tag{R29.12.4}
\]

The three extra edges are the original root unit arm and the two
outer edges. Since this is the original tree component,
|\(V_B|=e_B+1\) and it has \(e_B\) nonreceiver vertices.
All path subdivision vertices are counted, as are all paired-arm
vertices. No skeleton edge count replaces \(e_B\) or the final e(U).
